% Generated by tools/md2tex.py from drafts/paper_en.md. Edit the draft, then re-run the script.
\begin{abstract}
Subject-to-video generation renders a user-provided subject into a video described by text. Current systems are trained on paired reference--video data, which are scarce and biased: the reference is usually cropped from the target video, so training rewards pasting the reference into the frame rather than re-rendering the subject. We instead learn appearance from image pairs only, whose references deliberately differ from their targets, and borrow motion entirely from a frozen video diffusion transformer. Two mechanisms make static image supervision compatible with a video model. A signed motion adapter turns the backbone into a motion carrier whose scale $M$ moves continuously between a static regime, native motion, and amplified motion; running the carrier in its static regime during training lets static image targets shape appearance without teaching the model to be still. A paired-difference readout evaluates the reference branch with the real reference and with a learned null reference and adds only their difference to the carrier, so the reference never touches the motion computation and the reference-free branch for guidance is learned by construction. At inference, motion, text adherence, and reference adherence become three independent controls. On OpenS2V-Eval, with a copy-paste metric validated on synthetic anchors and gated by subject identity, the image-trained model copies the reference markedly less than the strongest video-trained baselines.
\end{abstract}
